% ============================================================================
% PAPER A · A pi-free generator of the lemniscatic field, with applications
%           to a closed-form quadratic for the fine-structure constant
% ============================================================================
% Target venue: Letters in Mathematical Physics (~10 pages).
% Strategy:     docs/theory/STRATEGY_PAPER_SPLIT_2026-04-30.md
% Created:      2026-05-01 evening (post physics-bridge crystallization).
% Citation targets:
%   - docs/theory/01_reference/SPEC_ALGEBRAIC_SPINE.md (Theorems 1, 2, 8, 9)
%   - docs/theory/03_derivations/THEOREM_HARMONIC_INVARIANT_TOWER.md (Theorem 8)
% Build:   cd dissemination/papers && pdflatex PAPER_A_PI_FREE_GENERATOR.tex
% Status:  DRAFT v2 (2026-05-02).
% ============================================================================

\documentclass[11pt,reqno]{amsart}

\usepackage[margin=1in]{geometry}
\usepackage{amsmath,amssymb,amsthm,mathtools}
\usepackage{booktabs}
\usepackage{hyperref}
\usepackage{xcolor}
\usepackage[numbers,sort&compress]{natbib}
\usepackage{enumitem}

\hypersetup{colorlinks=true,
  linkcolor=blue!50!black,
  citecolor=blue!50!black,
  urlcolor=blue!50!black}

\theoremstyle{plain}
\newtheorem{theorem}{Theorem}[section]
\newtheorem{lemma}[theorem]{Lemma}
\newtheorem{proposition}[theorem]{Proposition}
\newtheorem{corollary}[theorem]{Corollary}
\theoremstyle{definition}
\newtheorem{definition}[theorem]{Definition}
\theoremstyle{remark}
\newtheorem{remark}[theorem]{Remark}

\numberwithin{equation}{section}

\newcommand{\Gstar}{G_{\!*}}
\newcommand{\Aut}{\operatorname{Aut}}
\newcommand{\disc}{\operatorname{disc}}

\title[A $\pi$-free generator of the lemniscatic field]
{A $\pi$-free generator of the lemniscatic field,\\
 with applications to a closed-form quadratic\\
 for the fine-structure constant}

\author{Anonymous}

\author{C.\ Pacitanimulli}
\email{cpacitanimulli@gmail.com}

\date{\today}

\subjclass[2020]{Primary 11J81, 11G15;
                 Secondary 11M06, 33B15, 81V05}
\keywords{Lemniscatic constant; complex multiplication; class-number-one;
  Chowla--Selberg; algebraic independence; Chudnovsky's theorem;
  fine-structure constant.}

\begin{document}

\begin{abstract}
Let $\Gstar := \Gamma(1/4)/\Gamma(3/4)$. We record four theorems about
$\Gstar$ and a one-parameter polynomial tower built from it. The
quadratic $P(x) = x^{2} - 16\,\Gstar^{2}\,x + 16\,\Gstar^{3}$ has
closed-form roots $x_{\pm} = 8\Gstar^{2} \pm 4\Gstar\sqrt{4\Gstar^{2}
- \Gstar}$. Embedded in the level-indexed family $M_{k}(x) := x^{2} -
2^{k}\Gstar^{k-2}\,x + 2^{k}\Gstar^{k-1}$ for $k \geq 3$, the
normalised roots $y_{\pm} := x_{\pm}/\Gstar$ satisfy the harmonic
invariant $1/y_{+} + 1/y_{-} = 1$ at every level, and the level-$k$
discriminant correction $A_{k} := 2^{k-2}\Gstar^{k-3} - 1$ is
transcendental over $\mathbb{Q}$ for $k \geq 4$ (rational at $k=3$).
\emph{Conditional on Chudnovsky's 1976 algebraic-independence theorem
for $\pi$ and $\Gamma(1/4)$,} we show that $\mathbb{Q}(\Gstar)$ is a
$\pi$-free subfield of $\mathbb{Q}(\pi, \Gamma(1/4))$, in the precise
sense that $\mathbb{Q}(\Gstar) \cap \mathbb{Q}(\pi) = \mathbb{Q}$.
The numerical proximity $x_{+} \approx 137.036$ to the reciprocal
fine-structure constant $1/\alpha$ (1.26 ppm to CODATA 2022) is
recorded as a strongly motivated empirical conjecture, not a
derivation; the conditional closed form
$\alpha = 1/(2\Gstar) - \sqrt{4\Gstar - 1}/(4\,\Gstar^{3/2})$ follows
algebraically from $P(x_{+}) = 0$ together with the identification
$x_{+} = 1/\alpha$.
\end{abstract}

\maketitle

% ============================================================================
\section{Introduction}
\label{sec:intro}

The constant
\begin{equation}
\label{eq:gstar-def}
  \Gstar \;:=\; \frac{\Gamma(1/4)}{\Gamma(3/4)} \;=\; 2.95867512\ldots
\end{equation}
arises in three classical settings: as the value of the lemniscate
period for the curve $E\colon y^{2} = x^{3} - x$, as a closed-form
evaluation of an $L$-series for the imaginary quadratic field
$\mathbb{Q}(i)$ via Chowla--Selberg, and as the closed-form value of
Watson's three-dimensional body-centred-cubic integral
$W_{3} = \Gstar^{2}/(2\pi)$ \cite{watson1939, glasser-zucker}. The
present paper concerns four mathematical statements about $\Gstar$
and a closed-form algebraic application.

The mathematical content is summarised in
\S\S\ref{sec:gstar}--\ref{sec:qg}: an algebraic identity for $\Gstar$
(Theorem~\ref{thm:gstar}), a one-parameter polynomial family in
which $\Gstar$ enters with integer multipliers
(Theorems~\ref{thm:masterq} and \ref{thm:tower}), and a
field-theoretic statement that
$\mathbb{Q}(\Gstar)$ is $\pi$-free in the field
$\mathbb{Q}(\pi, \Gamma(1/4))$ \emph{conditional on the Chudnovsky
algebraic-independence theorem} (Theorem~\ref{thm:qgstar}).

The application appears in \S\ref{sec:closed-form}. The smallest
member of the polynomial family --- the quadratic
\begin{equation}
\label{eq:masterq-intro}
  P(x) \;=\; x^{2} - 16\,\Gstar^{2}\,x + 16\,\Gstar^{3}
\end{equation}
--- has two real roots, the larger of which has been observed
empirically to coincide with the reciprocal fine-structure constant
$1/\alpha$ to 1.26 parts per million on CODATA 2022 data. We take this
empirical observation as motivation but \emph{not} as derivation;
\S\ref{sec:closed-form} states the consequence as a corollary
conditional on the identification $x_{+} = 1/\alpha$, and reports the
algebraic closed form
\[
  \alpha \;=\; \frac{1}{2\Gstar} \;-\; \frac{\sqrt{4\Gstar - 1}}{4\,\Gstar^{3/2}}
\]
in that conditional sense.

The empirical identification $x_{+} = 1/\alpha$ is external to the
mathematical content of the paper.

A note on terminology. The constant $\Gstar = \Gamma(1/4)/\Gamma(3/4)
\approx 2.9587$ is sometimes called ``the lemniscate constant'' in
informal contexts; it is to be distinguished from the
Bernoulli--Gauss lemniscate constant $\varpi = \Gamma(1/4)^{2}/(2\sqrt{2\pi})
\approx 2.6221$. In this paper we reserve $\Gstar$ for the ratio
$\Gamma(1/4)/\Gamma(3/4)$ throughout, and write $\varpi$ explicitly
when the Bernoulli value is meant. The two are linked by
$\Gstar = 2\varpi/\sqrt{\pi}$.

\paragraph{Notation.} Throughout, $\mathbb{Q}$ denotes the rationals,
$\overline{\mathbb{Q}}$ their algebraic closure, and $\Gamma$ Euler's
gamma function. Numerical values quoted to many digits are computed
from \eqref{eq:gstar-def} and the elementary identity
$\Gamma(1/4)\,\Gamma(3/4) = \pi\sqrt{2}$.

% ============================================================================
\section{The constant $\Gstar$}
\label{sec:gstar}

\begin{theorem}[Algebraic identity for $\Gstar$]
\label{thm:gstar}
With $\Gstar$ defined by \eqref{eq:gstar-def}, the equivalent forms
\begin{equation}
\label{eq:gstar-forms}
  \Gstar \;=\; \frac{\Gamma(1/4)}{\Gamma(3/4)}
        \;=\; \frac{\Gamma(1/4)^{2}}{\pi\sqrt{2}}
        \;=\; \frac{2\,\varpi}{\sqrt{\pi}}
\end{equation}
hold, where $\varpi := \Gamma(1/4)^{2}/(2\sqrt{2\pi})$ is the
Bernoulli--Gauss lemniscate constant. Numerically
$\Gstar = 2.95867512\ldots$ and $16\,\Gstar^{2} = 140.060\ldots$.
\end{theorem}

\begin{proof}
The Euler reflection formula gives
$\Gamma(1/4)\,\Gamma(3/4) = \pi/\sin(\pi/4) = \pi\sqrt{2}$. Dividing
$\Gamma(1/4)^{2}$ through by this identity yields the second form;
the third follows from the definition of $\varpi$.
\end{proof}

The two further classical evaluations of $\Gstar$ --- via the
Chowla--Selberg formula applied to the elliptic curve
$E\colon y^{2} = x^{3} - x$, and via Watson's $W_{3}$ identity
\cite{watson1939} --- are consequences of $\eqref{eq:gstar-forms}$
together with standard facts about the lemniscatic period
\cite{chowla-selberg, glasser-zucker}; we cite them but do not
re-prove them here. They serve to anchor $\Gstar$ as a constant of
established mathematical literature, not a contrivance for this
paper.

% ============================================================================
\section{The master quadratic}
\label{sec:masterq}

\begin{theorem}[Master quadratic and roots]
\label{thm:masterq}
Let $P(x) := x^{2} - 16\,\Gstar^{2}\,x + 16\,\Gstar^{3}$. The
discriminant
\begin{equation}
\label{eq:disc}
  \disc(P) \;=\; 256\,\Gstar^{4} - 64\,\Gstar^{3}
            \;=\; 64\,\Gstar^{3}(4\,\Gstar - 1)
\end{equation}
is positive since $\Gstar > 1/4$. Hence $P$ has two distinct real
roots
\begin{equation}
\label{eq:roots}
  x_{\pm} \;=\; 8\,\Gstar^{2} \;\pm\; 4\,\Gstar\,\sqrt{4\,\Gstar^{2} - \Gstar}.
\end{equation}
Numerically $x_{+} = 137.036171\ldots$ and $x_{-} = 3.023964\ldots$.
\end{theorem}

\begin{proof}
Direct application of the quadratic formula and computation of the
discriminant. Positivity of $\disc(P)$ uses only $\Gstar > 1/4$,
which is immediate from \eqref{eq:gstar-forms} and
$\Gamma(1/4) > \Gamma(3/4) > 0$.
\end{proof}

The polynomial $P$ has integer-and-$\Gstar$-power coefficients, with
both coefficients sharing the integer factor 16. We do not assign a
mathematical meaning to the value 16 in this paper beyond ``the
smallest integer for which the resulting two real roots are positive
and admit a clean closed form''. (The factor 16 also appears as
$|\Aut(E)|^{2}$ for $E\colon y^{2} = x^{3}-x$, which has automorphism
group of order 4 over $\overline{\mathbb{Q}}$; this connection is not
used in any argument below and is mentioned only to fix the reader's
intuition.)

% ============================================================================
\section{The harmonic-invariant tower}
\label{sec:tower}

The master quadratic is the level-$k=4$ instance of a one-parameter
family.

\begin{definition}[The level-indexed family]
For each integer $k \geq 3$, define
\begin{equation}
\label{eq:tower-def}
  M_{k}(x) \;:=\; x^{2} \;-\; 2^{k}\,\Gstar^{k-2}\,x \;+\; 2^{k}\,\Gstar^{k-1}.
\end{equation}
\end{definition}

The $k=4$ instance is the master quadratic of
Theorem~\ref{thm:masterq}, since $2^{4} = 16$. Each $M_{k}$ has two
real roots $x_{\pm}^{(k)}$; we write $y_{\pm}^{(k)} :=
x_{\pm}^{(k)}/\Gstar$ for the normalised roots.

\begin{theorem}[Harmonic invariant; level-$k$ discriminant transcendence]
\label{thm:tower}
For every $k \geq 3$:
\begin{enumerate}[label=\textup{(\roman*)}, leftmargin=*]
\item \emph{Harmonic invariant.} $1/y_{+}^{(k)} + 1/y_{-}^{(k)} = 1$.
\item \emph{Discriminant factorisation.}
$\disc(M_{k}) = 2^{k+2}\,\Gstar^{k-1}\,A_{k}$, where
\begin{equation}
\label{eq:Ak}
  A_{k} \;:=\; 2^{k-2}\,\Gstar^{k-3} - 1.
\end{equation}
\item \emph{Discriminant transcendence.} $A_{3} = 1 \in \mathbb{Q}$,
and $A_{k} \in \mathbb{R} \setminus \overline{\mathbb{Q}}$ for every
$k \geq 4$ \emph{conditional on the Schneider--Chudnovsky theorem
\cite{chudnovsky1976}.}
\end{enumerate}
\end{theorem}

\begin{proof}
\textit{(i)} By Vieta's relations applied to $M_{k}$,
$x_{+}^{(k)} + x_{-}^{(k)} = 2^{k}\,\Gstar^{k-2}$ and
$x_{+}^{(k)}\,x_{-}^{(k)} = 2^{k}\,\Gstar^{k-1}$. Hence
\[
  \frac{1}{x_{+}^{(k)}} + \frac{1}{x_{-}^{(k)}}
  \;=\; \frac{x_{+}^{(k)} + x_{-}^{(k)}}{x_{+}^{(k)}\,x_{-}^{(k)}}
  \;=\; \frac{2^{k}\,\Gstar^{k-2}}{2^{k}\,\Gstar^{k-1}}
  \;=\; \frac{1}{\Gstar}.
\]
Multiplying through by $\Gstar$ gives
$1/y_{+}^{(k)} + 1/y_{-}^{(k)} = 1$.

\textit{(ii)} Direct expansion:
\[
  \disc(M_{k}) \;=\; (2^{k}\Gstar^{k-2})^{2} - 4 \cdot 2^{k}\Gstar^{k-1}
              \;=\; 2^{2k}\Gstar^{2k-4} - 2^{k+2}\Gstar^{k-1}
              \;=\; 2^{k+2}\Gstar^{k-1}\bigl(2^{k-2}\Gstar^{k-3} - 1\bigr).
\]

\textit{(iii)} At $k=3$, $A_{3} = 2^{1}\Gstar^{0} - 1 = 1$. For
$k \geq 4$, $A_{k} = 2^{k-2}\Gstar^{k-3} - 1$ is a polynomial in
$\Gstar$ with rational coefficients, of degree $k - 3 \geq 1$, with
non-zero leading coefficient $2^{k-2}$. By the Chudnovsky
theorem \cite{chudnovsky1976}, $\pi$ and $\Gamma(1/4)$ are
\emph{algebraically independent} over $\mathbb{Q}$ --- a statement
strictly stronger than the transcendence of $\Gamma(1/4)$ alone
(Schneider, 1948), and the one needed here; together with the algebraic
relation $\Gstar\cdot\pi\sqrt{2} = \Gamma(1/4)^{2}$ from
\eqref{eq:gstar-forms}, this forces $\Gstar$ to be transcendental over
$\mathbb{Q}$. A
non-rational polynomial in a transcendental over $\mathbb{Q}$ with
rational coefficients takes only transcendental values
\cite[\S 1.4]{waldschmidt2000}. Hence $A_{k} \notin \overline{\mathbb{Q}}$.
\end{proof}

\begin{remark}
The $k=3$ instance $M_{3}(x) = x^{2} - 8\,x + 8\,\Gstar$ has roots
$x_{\pm}^{(3)} = 4 \pm \sqrt{16 - 8\Gstar}$. Since
$16 - 8\Gstar = 8(2 - \Gstar) < 0$ for our $\Gstar \approx 2.96$, the
$k=3$ roots are complex conjugates, and the harmonic invariant
$1/y_{+} + 1/y_{-} = 1$ holds in the complex sense. The first
\emph{real-rooted} level is $k = 4$, the master quadratic.
\end{remark}

The harmonic invariant $1/y_{+} + 1/y_{-} = 1$ characterises the
level-indexed family in the following sense. Among one-parameter
families of monic quadratics
$\{Q(x) = x^{2} - bx + c\}$ with $b, c \in \mathbb{R}_{>0}$, the
identity $1/y_{+} + 1/y_{-} = 1$ (for $y = x/\Gstar$) holds precisely
when $c = \Gstar \cdot b$. The family $\{M_{k}\}_{k \geq 3}$ is a
particular sub-family of this constraint, parameterised by the
integer multiplier $2^{k}$.

% ============================================================================
\section{$\mathbb{Q}(\Gstar)$ is $\pi$-free in $\mathbb{Q}(\pi, \Gamma(1/4))$}
\label{sec:qg}

\begin{theorem}[Conditional on Chudnovsky 1976]
\label{thm:qgstar}
\emph{Conditional on the Chudnovsky algebraic-independence theorem
for $\pi$ and $\Gamma(1/4)$ \cite{chudnovsky1976},}
$\mathbb{Q}(\Gstar)$ is a subfield of $\mathbb{Q}(\pi, \Gamma(1/4))$
satisfying
\begin{equation}
\label{eq:qg-trivial-int}
  \mathbb{Q}(\Gstar) \;\cap\; \mathbb{Q}(\pi) \;=\; \mathbb{Q}.
\end{equation}
That is, $\Gstar$ generates a subfield of
$\mathbb{Q}(\pi, \Gamma(1/4))$ whose intersection with
$\mathbb{Q}(\pi)$ is trivial.
\end{theorem}

\begin{proof}
\emph{Containment.} By \eqref{eq:gstar-forms},
$\Gstar = \Gamma(1/4)^{2}/(\pi\sqrt{2})$, a rational function in
$\{\Gamma(1/4), \pi\}$ over $\mathbb{Q}(\sqrt{2})$, hence over
$\mathbb{Q}$ after adjoining the algebraic element $\sqrt{2}$. Every
element of $\mathbb{Q}(\Gstar)$ is therefore a rational function in
$\{\Gamma(1/4), \pi\}$, i.e.\ lies in $\mathbb{Q}(\pi, \Gamma(1/4))$
after the algebraic adjunction (which we absorb into
$\overline{\mathbb{Q}}$ for the purposes of intersection with
$\mathbb{Q}(\pi)$).

\emph{Trivial intersection.} Suppose $\beta \in \mathbb{Q}(\Gstar) \cap
\mathbb{Q}(\pi)$. Then $\beta = p(\Gstar)/q(\Gstar)$ for some
$p, q \in \mathbb{Q}[T]$ with $q(\Gstar) \neq 0$, and simultaneously
$\beta = f(\pi)/g(\pi)$ for some $f, g \in \mathbb{Q}[T]$ with
$g(\pi) \neq 0$. Cross-multiplying yields
\begin{equation}
\label{eq:cross-mult}
  p(\Gstar)\,g(\pi) \;=\; q(\Gstar)\,f(\pi),
\end{equation}
a polynomial relation in $\Gstar$ and $\pi$ with rational coefficients.

We claim $\Gstar$ and $\pi$ are algebraically independent over
$\mathbb{Q}$. By Chudnovsky's 1976 theorem,
$\{\pi, \Gamma(1/4)\}$ is algebraically independent over $\mathbb{Q}$.
Suppose for contradiction there is a non-trivial
$P \in \mathbb{Q}[X, Y]$ with $P(\Gstar, \pi) = 0$. Substituting
$\Gstar = \Gamma(1/4)^{2}/(\pi\sqrt{2})$ into $P$ and clearing
denominators (multiplying by an appropriate power of $\pi$) yields a
polynomial $\widetilde{P} \in \mathbb{Q}(\sqrt{2})[X, Y]$ with
$\widetilde{P}(\Gamma(1/4), \pi) = 0$. Since $\sqrt{2}$ is algebraic
over $\mathbb{Q}$, $\widetilde{P}$ has coefficients in
$\overline{\mathbb{Q}}$, contradicting the algebraic independence of
$\{\pi, \Gamma(1/4)\}$ over $\mathbb{Q}$ (which holds equally over
$\overline{\mathbb{Q}}$ since the property is preserved under
algebraic extensions). Hence $\Gstar$ and $\pi$ are algebraically
independent over $\mathbb{Q}$.

Returning to \eqref{eq:cross-mult}: the only polynomial relation
$P(\Gstar, \pi) = 0$ consistent with algebraic independence is the
trivial one in which both sides reduce identically. This forces
$\beta$ to be a constant, hence $\beta \in \mathbb{Q}$. \qedhere
\end{proof}

\begin{remark}[On the conditionality]
The theorem rests on Chudnovsky's 1976 result \cite{chudnovsky1976}, a
standard result of contemporary transcendence theory consolidated in
\cite[\S 1.4]{waldschmidt2000}. ``Conditional'' here means
``depends on this established theorem'', not ``depends on a
conjecture''. The conditionality is recorded in the abstract, the
theorem statement, and the conclusion (\S\ref{sec:discussion}) in
keeping with the convention that any result whose proof appeals to a
named external theorem in transcendence theory should bear the cite
visibly.
\end{remark}

\begin{remark}[What the theorem does not claim]
Theorem~\ref{thm:qgstar} does \emph{not} claim that
$\mathbb{Q}(\Gstar)$ is a maximal $\pi$-free subfield of
$\mathbb{Q}(\pi, \Gamma(1/4))$ in any extremal sense. Larger
$\pi$-free subfields may exist; the result only establishes
$\mathbb{Q}$-triviality of the intersection with $\mathbb{Q}(\pi)$.
\end{remark}

% ============================================================================
\section{Closed form for $\alpha$ conditional on the empirical match}
\label{sec:closed-form}

The numerical value $x_{+} = 137.036171\ldots$ from
Theorem~\ref{thm:masterq} matches the reciprocal fine-structure
constant $1/\alpha$ to 1.26 parts per million on CODATA 2022 data.
Specifically, the CODATA 2022 recommended value is
$\alpha^{-1} = 137.035999177(21)$ with quoted relative uncertainty
$\approx 1.5 \times 10^{-10}$ \cite{codata2022}. The discrepancy
\[
  |x_{+} - \alpha^{-1}| / \alpha^{-1} \;=\; 1.26 \times 10^{-6}
\]
exceeds CODATA precision by four orders of magnitude; the match is
empirical, not exact.

We do not derive $\alpha = 1/x_{+}$ from any of the mathematical
statements above. The empirical proximity is recorded as a strongly
motivated conjecture; in this paper it serves only as motivation.

\begin{corollary}[Closed form, conditional]
\label{cor:closed-form}
If the empirical identification $x_{+} = 1/\alpha$ holds exactly,
then
\begin{equation}
\label{eq:closed-form-alpha}
  \alpha \;=\; \frac{1}{2\,\Gstar} \;-\; \frac{\sqrt{4\,\Gstar - 1}}{4\,\Gstar^{3/2}}.
\end{equation}
\end{corollary}

\begin{proof}
By Theorem~\ref{thm:masterq}, $x_{+} = 8\,\Gstar^{2} +
4\,\Gstar\sqrt{4\,\Gstar^{2} - \Gstar}$. Substituting
$\alpha = 1/x_{+}$ and rationalising gives
\[
  \alpha \;=\; \frac{1}{8\,\Gstar^{2} + 4\,\Gstar\,\sqrt{4\,\Gstar^{2} - \Gstar}}
        \;=\; \frac{8\,\Gstar^{2} - 4\,\Gstar\,\sqrt{4\,\Gstar^{2} - \Gstar}}
                   {(8\Gstar^{2})^{2} - 16\,\Gstar^{2}(4\,\Gstar^{2} - \Gstar)}
        \;=\; \frac{8\,\Gstar^{2} - 4\,\Gstar\,\sqrt{4\,\Gstar^{2} - \Gstar}}
                   {16\,\Gstar^{3}}.
\]
Splitting the fraction and simplifying $4\,\Gstar^{2} - \Gstar
= \Gstar(4\,\Gstar - 1)$ gives \eqref{eq:closed-form-alpha}.
\end{proof}

\begin{remark}
Corollary~\ref{cor:closed-form} is purely algebraic given the
identification. It does \emph{not} establish that $\alpha = 1/x_{+}$;
that identification is the conjecture, and the corollary is its
consequence. The role of \eqref{eq:closed-form-alpha} in the paper is
to make explicit that \emph{if} the conjecture holds, then $\alpha$
admits a closed form in $\Gstar$ alone --- equivalently, in
$\Gamma(1/4)$ and $\pi$ via \eqref{eq:gstar-forms}.
\end{remark}

% ============================================================================
\section{Discussion}
\label{sec:discussion}

The four mathematical statements (Theorems \ref{thm:gstar},
\ref{thm:masterq}, \ref{thm:tower}, \ref{thm:qgstar}) hold
unconditionally except for the third part of Theorem~\ref{thm:tower}
(transcendence of $A_{k}$ for $k \geq 4$) and Theorem~\ref{thm:qgstar}
(triviality of $\mathbb{Q}(\Gstar) \cap \mathbb{Q}(\pi)$), both of
which are conditional on Chudnovsky's 1976 algebraic-independence
theorem \cite{chudnovsky1976}. The conditionality is named in each
statement, repeated in the abstract, and re-iterated here: the four
theorems together form an unconditional core (Theorems \ref{thm:gstar},
\ref{thm:masterq}, \ref{thm:tower}(i)--(ii)) and a Chudnovsky-conditional
extension (Theorems \ref{thm:tower}(iii) and \ref{thm:qgstar}).

The closed-form expression for $\alpha$ in
Corollary~\ref{cor:closed-form} is conditional on the identification
$x_{+} = 1/\alpha$, which is empirical and at present unmotivated by
any derivation we know of. Its mathematical content is the algebraic
inversion of $P$ at the larger root; its physical content is exactly
what one chooses to make of the 1.26-ppm empirical match.

We close with two observations on possible directions.

\paragraph{(1) Are there larger $\pi$-free subfields of $\mathbb{Q}(\pi, \Gamma(1/4))$?}
Theorem~\ref{thm:qgstar} establishes that $\mathbb{Q}(\Gstar)$ is one
$\pi$-free subfield, with trivial intersection in the sense of
\eqref{eq:qg-trivial-int}. It does not characterise all such
subfields, nor does it claim maximality. A natural question is
whether $\mathbb{Q}(\Gstar)$ is maximal among $\pi$-free subfields
under appropriate conditions; we leave this open.

\paragraph{(2) Tower levels $k > 4$.}
The harmonic invariant of Theorem~\ref{thm:tower}(i) holds at every
level $k \geq 3$. Levels $k = 4$ and $k = 5$ produce real roots; we
have not investigated whether any higher level produces roots whose
numerical proximity to other dimensionless constants of physics is
similarly tight. Should such a level be identified, the corresponding
discriminant correction $A_{k}$ would by Theorem~\ref{thm:tower}(iii)
be transcendental for $k \geq 4$, and the analogous closed form would
follow from the same Vieta inversion as Corollary~\ref{cor:closed-form}.

% ============================================================================
\section*{Acknowledgements}

The Chowla--Selberg evaluation underlying $\Gstar$, Watson's $W_{3}$
identity, and Chudnovsky's algebraic-independence theorem are
classical results of analytic number theory and transcendence theory;
we have used them freely.

% ============================================================================
\begin{thebibliography}{99}

\bibitem{chowla-selberg}
S.\ Chowla and A.\ Selberg, \textit{On Epstein's zeta-function}, J.\
Reine Angew.\ Math.\ \textbf{227} (1967), 86--110.

\bibitem{chudnovsky1976}
D.\ V.\ Chudnovsky, \textit{Algebraic independence of values of
exponential and hypergeometric functions}, in: \textit{Number Theory}
(Proc.\ Conf., Carbondale, IL, 1979), Lecture Notes in Math.\
\textbf{751}, Springer, Berlin, 1979, 45--69; result originally
announced 1976.

\bibitem{codata2022}
E.\ Tiesinga, P.\ J.\ Mohr, D.\ B.\ Newell, and B.\ N.\ Taylor,
\textit{CODATA recommended values of the fundamental physical
constants: 2022}, Rev.\ Mod.\ Phys.\ \textbf{96} (2024), 025002.

\bibitem{glasser-zucker}
M.\ L.\ Glasser and I.\ J.\ Zucker, \textit{Lattice sums}, in:
\textit{Theoretical Chemistry: Advances and Perspectives}, vol.\ 5,
Academic Press, 1980, 67--139.

\bibitem{waldschmidt2000}
M.\ Waldschmidt, \textit{Diophantine Approximation on Linear
Algebraic Groups: Transcendence Properties of the Exponential
Function in Several Variables}, Grundlehren der mathematischen
Wissenschaften \textbf{326}, Springer, Berlin, 2000.

\bibitem{watson1939}
G.\ N.\ Watson, \textit{Three triple integrals}, Quart.\ J.\ Math.\
Oxford Ser.\ \textbf{10} (1939), 266--276.

\end{thebibliography}

\end{document}
